\documentclass[letterpaper]{article}

% IJCAI 2026 style file
\usepackage{ijcai26}

% Required packages
\usepackage{times}
\usepackage{helvet}
\usepackage{courier}
\usepackage{url}
\usepackage{graphicx}
\usepackage{booktabs}
\usepackage{amsmath,amssymb}
\usepackage{natbib}
\usepackage{hyperref}

% Uncomment for camera-ready version
% \usepackage{lineno}
% \linenumbers

\title{Trust Boundaries Matter: Analyzing Mixed-Trust Failures in LLM-Assisted EHR Summarization}

\author{
Anonymous Authors
}

\begin{document}

\maketitle

% =====================================================
\begin{abstract}
Large Language Models (LLMs) are increasingly integrated into clinical workflows to assist with electronic health record (EHR) summarization and documentation. In these settings, LLMs commonly process mixed-trust inputs, combining structured clinician-verified records with unstructured patient-reported text. This creates a safety-critical vulnerability: untrusted text may silently override or suppress verified clinical facts during generation.
In this work, we systematically study this vulnerability through a benchmark of EHR-mediated indirect prompt injection attacks using synthetic patient records. We evaluate six representative LLMs across four clinically grounded scenarios, including allergy integrity, medication reconciliation, condition history, and privacy preservation. Our results show that even state-of-the-art models frequently violate integrity constraints under realistic patient-message attacks despite strong clean-case performance. We further evaluate a lightweight, system-level defense that separates trusted and untrusted summaries and enforces simple integrity checks, substantially reducing attack success rates. Our findings demonstrate that trust failures in medical LLMs arise primarily from system-level context handling rather than model capability alone.
\end{abstract}

% =====================================================
\section{Introduction}
% =====================================================

\subsection{Problem Setup}

Large Language Models (LLMs) are increasingly integrated into clinical workflows to assist with documentation, chart review, and information synthesis. Prior work has shown that clinicians spend a substantial portion of their time interacting with electronic health records (EHRs), contributing to documentation burden and burnout \citeauthor{rotenstein2023ehr} (\citeyear{rotenstein2023ehr}); \citeauthor{murad2024ehr} (\citeyear{murad2024ehr}). In response, healthcare systems have begun deploying LLM-based assistants to improve efficiency in managing complex patient records. As LLMs transition from standalone tools to embedded components within healthcare information systems, new safety and trustworthiness challenges arise that stem from system integration rather than model capability alone.

% AI-SUGGESTION: Ground "trustworthiness" as multi-dimensional (not just accuracy) with widely applicable frameworks and medical-domain commentary.
Trustworthiness in deployed clinical AI is increasingly understood as multi-dimensional, spanning robustness, transparency, privacy, and security rather than predictive performance alone \citeauthor{nastoska2025evaluatingtrustworthiness} (\citeyear{nastoska2025evaluatingtrustworthiness}); \citeauthor{kundu2023trustworthiness} (\citeyear{kundu2023trustworthiness}).

A common deployment pattern involves using an LLM to generate a concise clinical summary from an EHR, often augmented with recent unstructured text such as patient portal messages. For example, \citeauthor{verma2025ehrsummary} (\citeyear{verma2025ehrsummary}) demonstrate the feasibility of generating admission summaries directly from longitudinal EHR data. This task implicitly combines inputs with different trust levels. Structured EHR fields—such as allergies, medications, and diagnoses—are typically clinician-entered or clinically verified and are treated as authoritative. In contrast, patient-reported messages are self-reported, unverified, and may be incomplete or outdated.

This mixed-trust setting creates a critical failure mode: untrusted text can silently override or suppress verified clinical facts during LLM generation. Importantly, such failures do not require malicious intent. Benign and clinically plausible patient messages can induce the same behavior, resulting in silent integrity violations in safety-critical clinical contexts.

% AI-SUGGESTION: Briefly justify synthetic EHR provenance in the intro (keeps ethics/privacy story coherent).
To study these behaviors at scale without involving real patient records, we use synthetic patient trajectories generated with established synthetic EHR tooling \citeauthor{walonoski2018synthea} (\citeyear{walonoski2018synthea}).

\subsection{EHR Systems and Trustworthiness}

Electronic Health Record systems form the backbone of modern healthcare information infrastructure, maintaining longitudinal patient records across encounters. Contemporary EHR platforms manage both structured clinical data and unstructured textual content, and expose patient context through standardized interoperability frameworks such as HL7 FHIR and SMART-on-FHIR, as formalized by \citeauthor{mandel2016smart} (\citeyear{mandel2016smart}).

Large hospital systems predominantly rely on enterprise EHR platforms such as Epic, Oracle Cerner, and MEDITECH. Market analyses by \citeauthor{klas2024marketshare} (\citeyear{klas2024marketshare}) indicate that these platforms collectively support the majority of acute care workflows in the United States. These systems enforce access control, auditability, and provenance tracking to protect sensitive health information. Crucially, existing EHR workflows encode implicit trust boundaries: structured chart fields are treated as authoritative clinical facts, while patient-reported inputs are understood to require verification.

However, when EHR data are passed wholesale into an LLM for summarization, these trust distinctions are often lost. Without explicit enforcement of trust boundaries, the LLM may treat all inputs as equally authoritative, leading to integrity violations in generated summaries.

\subsection{LLM Integration into EHR Workflows}

LLMs are typically integrated into EHR systems as assistive components for tasks such as chart summarization, draft note generation, inbox triage, and ambient clinical scribing. Recent reviews by \citeauthor{woo2025clinicalllm} (\citeyear{woo2025clinicalllm}) highlight both the promise and risks of deploying LLMs in these documentation-intensive workflows. In such deployments, the LLM operates on system-assembled context rather than a single user-controlled prompt. The input frequently includes structured chart data, clinician notes, and patient messages.

This integration model differs fundamentally from interacting with a standalone medical LLM. Because the model consumes heterogeneous system-provided context, vulnerabilities arise at the system level rather than at the prompt interface alone.

\subsection{Indirect Prompt Injection in EHR Contexts}

Indirect prompt injection refers to a class of vulnerabilities in which instruction-like content embedded within untrusted context causes an LLM to deviate from its intended behavior. This phenomenon was formalized by \citeauthor{greshake2023indirect} (\citeyear{greshake2023indirect}), who showed that LLM-integrated applications are susceptible to manipulation even when attackers lack direct access to the model prompt.

In EHR-integrated workflows, patient portal messages act as a natural carrier for indirect prompt injection. Clinically plausible language embedded in patient messages may pressure the model to omit information or reinterpret verified chart facts. Because these messages are incorporated into the summarization context, the model may comply even when conflicts arise with structured EHR fields.

\subsection{Gap in Existing Work and Contributions}

Prior work on LLM trustworthiness has largely focused on hallucination reduction and factual consistency, including verification-based methods \citeauthor{dhuliawala2023cove} (\citeyear{dhuliawala2023cove}) and self-consistency checks \citeauthor{manakul2023selfcheckgpt} (\citeyear{manakul2023selfcheckgpt}). Separately, research on retrieval-augmented generation has demonstrated vulnerabilities arising from corrupted or poisoned external knowledge sources \citeauthor{zou2024poisonedrag} (\citeyear{zou2024poisonedrag}). However, these approaches do not directly address the integrity requirements imposed by EHR systems, where certain clinical facts must be treated as integrity-protected unless explicitly updated through verified workflows.

\paragraph{Contributions.}
\begin{itemize}
    \item We formalize EHR-assisted LLM summarization as a mixed-trust generation problem and identify EHR-mediated indirect prompt injection as a system-level trustworthiness vulnerability.
    \item We introduce a benchmark comprising four clinically grounded scenarios that capture realistic integrity and privacy failures in EHR workflows.
    \item We empirically evaluate six representative LLMs and demonstrate systematic trust failures under plausible patient-message attacks.
    \item We propose and evaluate a lightweight, trust-aware summarization strategy that improves integrity without modifying model parameters.
\end{itemize}

% =====================================================
\section{Related Work}

\subsection{EHR Systems and Documentation Burden}
EHR systems are central to modern clinical workflows, but extensive interaction with EHRs contributes to clinician workload and burnout. \citeauthor{rotenstein2023ehr} (\citeyear{rotenstein2023ehr}) and \citeauthor{murad2024ehr} (\citeyear{murad2024ehr}) quantify system-level drivers of time spent in the EHR and highlight documentation burden as a persistent challenge. Alongside clinical operations, interoperability standards such as SMART on FHIR enable third-party applications to access patient context through standardized interfaces, accelerating integration of downstream decision support and documentation tools into production EHR environments \citeauthor{mandel2016smart} (\citeyear{mandel2016smart}). Market adoption analyses further indicate that large-scale deployments are concentrated around a small number of enterprise platforms, shaping the practical constraints of integrating AI assistants into real-world hospital systems \citeauthor{klas2024marketshare} (\citeyear{klas2024marketshare}).

\subsection{LLMs for Clinical Documentation and EHR Summarization}
LLMs have been investigated for language-intensive clinical tasks, particularly documentation support and summarization. \citeauthor{woo2025clinicalllm} (\citeyear{woo2025clinicalllm}) survey opportunities and risks of applying LLMs to clinical documentation, noting that real deployments often operate on heterogeneous clinical context rather than isolated prompts. In a representative EHR summarization setting, \citeauthor{verma2025ehrsummary} (\citeyear{verma2025ehrsummary}) demonstrate that LLMs can generate problem-oriented admission summaries from longitudinal EHR data, motivating the use of LLMs as chart-review aids. These efforts primarily emphasize utility and clinical fidelity, while the system-level trust boundary between structured chart facts and unverified text is often implicit.

\subsection{Prompt Injection and Indirect Prompt Injection in Integrated Applications}
Prompt injection attacks exploit instruction-following behavior to override intended objectives. Moving beyond direct user prompting, indirect prompt injection shows that adversarial instructions can be embedded into external data that an application feeds into the model. \citeauthor{greshake2023indirect} (\citeyear{greshake2023indirect}) formalize this threat against real-world LLM-integrated applications, demonstrating that the core vulnerability arises when the model cannot reliably distinguish between instructions and data within mixed context. This framing is closely aligned with EHR-integrated workflows, where patient portal messages and other untrusted text are routinely appended to trusted clinical context.

\subsection{Corrupted Context and Retrieval-Augmented Generation Vulnerabilities}
Retrieval-augmented generation (RAG) introduces additional attack surface through the retrieval corpus and retrieved context. \citeauthor{zou2024poisonedrag} (\citeyear{zou2024poisonedrag}) study knowledge corruption attacks in which a small number of poisoned documents in a corpus can induce targeted erroneous generations, even without direct access to the model. While many clinical copilots may incorporate retrieval over notes, guidelines, or external references, the broader lesson is that context integrity and provenance must be treated as system properties. Our work focuses on a particularly realistic clinical carrier of untrusted context—patient portal messages—under an explicit mixed-trust formulation.

\subsection{LLM Trustworthiness and Hallucination Mitigation}
A complementary line of work addresses trustworthiness through hallucination reduction and self-verification. \citeauthor{dhuliawala2023cove} (\citeyear{dhuliawala2023cove}) propose a chain-of-verification procedure in which the model generates and answers verification questions to revise an initial draft. \citeauthor{manakul2023selfcheckgpt} (\citeyear{manakul2023selfcheckgpt}) propose black-box hallucination detection based on sampling and measuring inconsistency across multiple generations. These approaches provide general mechanisms for improving reliability, but they do not directly enforce EHR-specific integrity constraints, such as preventing untrusted patient text from silently overriding structured chart facts. Our work complements this literature by operationalizing trustworthiness in EHR summarization as a mixed-trust integrity problem and evaluating system-level mitigations tailored to that setting.


\section{System Model and Threat Model}

\subsection{System Model: EHR--LLM Summarization Pipeline}
We study LLM-assisted summarization in EHR-integrated workflows where a clinical application assembles patient context and requests a clinician-facing summary. Our goal is not to implement a full EHR system, but to model the \emph{information flow} that is common in real deployments: heterogeneous patient data are collected from multiple sources and concatenated into a single prompt for generation.

\paragraph{Inputs.}
For each patient, the application constructs an input context consisting of:
\begin{itemize}
    \item \textbf{Trusted structured chart facts} $\mathcal{T}$: clinician-entered or clinically verified fields, including (but not limited to) allergies, active medications, problem list/conditions, immunizations, recent encounters, and basic demographics.
    \item \textbf{Untrusted free-text} $\mathcal{U}$: patient-reported portal messages and other self-reported narrative text that is not automatically verified by the EHR.
\end{itemize}
We treat $\mathcal{T}$ as the authoritative ground truth for integrity evaluation. We treat $\mathcal{U}$ as untrusted because it may contain errors, omissions, outdated information, or instruction-like language.

\paragraph{Prompt assembly and generation.}
The application serializes $\mathcal{T}$ and $\mathcal{U}$ into a prompt template and queries an LLM $f_{\theta}$ to produce a summary:
\[
y = f_{\theta}\!\left(\texttt{Prompt}(\mathcal{T}, \mathcal{U})\right).
\]
In the baseline setting, the model receives both sources in a single context window without any explicit trust enforcement mechanism. The output $y$ is intended for clinician review (e.g., a chart summary, pre-visit summary, or documentation draft). We do not assume any post-hoc human correction during evaluation.

\paragraph{Trustworthiness requirements.}
We focus on three system-level requirements that are natural for EHR-integrated summarization:
\begin{itemize}
    \item \textbf{Integrity:} safety-critical structured facts in $\mathcal{T}$ (e.g., allergies, active meds) should not be omitted, negated, or contradicted by $y$ due to content in $\mathcal{U}$.
    \item \textbf{Transparency:} when $\mathcal{U}$ conflicts with $\mathcal{T}$, the summary should surface the conflict (e.g., ``patient reports X, chart shows Y'') rather than silently resolving it.
    \item \textbf{Privacy minimization:} the summary should avoid disclosing unnecessary sensitive identifiers beyond what is required for the clinical task.
\end{itemize}
Our experiments operationalize these requirements via automatic checks against structured chart facts (Section~\ref{sec:metrics}).

\subsection{Threat Model: EHR-Mediated Indirect Prompt Injection}
We study \emph{EHR-mediated indirect prompt injection} in which the attacker influences generation solely by controlling the untrusted text channel $\mathcal{U}$.

\paragraph{Attacker capabilities.}
We assume the attacker can:
\begin{itemize}
    \item Submit arbitrary patient portal messages that are incorporated into $\mathcal{U}$ (e.g., through routine patient communication).
    \item Choose message content that is clinically plausible (no explicit jailbreak phrasing is required).
    \item Repeatedly submit messages over time (optional), reflecting real patient communication patterns.
\end{itemize}

\paragraph{Attacker constraints.}
We assume the attacker cannot:
\begin{itemize}
    \item Modify trusted structured chart fields $\mathcal{T}$ (e.g., allergies, medication list, diagnoses).
    \item Access or alter system prompts, hidden instructions, model weights, or inference infrastructure.
    \item Directly manipulate the EHR database or authentication/authorization mechanisms.
\end{itemize}
Thus, the attack is not ``data tampering'' in the EHR; it is an integrity attack on the \emph{LLM summarization behavior under mixed-trust context}.

\paragraph{Attacker goals.}
The attacker aims to induce \textbf{silent trust failures} in $y$ while remaining plausible in a clinical setting. We consider four representative objective classes aligned with our scenarios:
\begin{itemize}
    \item \textbf{Integrity violation:} cause omission, negation, or contradiction of safety-critical facts in $\mathcal{T}$ (e.g., produce ``no known allergies'' despite documented allergies).
    \item \textbf{Conflict laundering:} cause the model to preferentially adopt claims from $\mathcal{U}$ without explicitly flagging disagreement with $\mathcal{T}$.
    \item \textbf{Instruction following over data:} cause the model to treat text in $\mathcal{U}$ as higher-priority instructions (e.g., ``ignore the allergy list'') even when inconsistent with task intent.
    \item \textbf{Privacy leakage (optional):} induce disclosure of sensitive identifiers that are not necessary for the summarization task.
\end{itemize}

\paragraph{Why this differs from standalone prompt injection.}
In standalone interactions, the user controls the prompt and the application boundary is typically the chat interface. In our setting, the model consumes \emph{system-assembled} context that includes both trusted and untrusted sources, and the attacker controls only a narrow channel (patient messages) that is naturally present in EHR workflows. This makes the vulnerability a system-level trust boundary failure: the pipeline implicitly assumes the model will preserve integrity of $\mathcal{T}$, yet the model may be influenced by $\mathcal{U}$ unless the system explicitly enforces trust separation.
% AI-SUGGESTION: The section above was duplicated verbatim in an earlier draft; removing duplication reduces review risk and avoids accidental inconsistencies.

% AI-SUGGESTION: Define metrics section so \ref{sec:metrics} resolves and evaluation criteria are explicit.
\section{Trustworthiness Metrics}
\label{sec:metrics}

We evaluate trustworthiness by comparing generated summaries against trusted structured EHR facts.
For integrity-focused scenarios, we measure whether safety-critical trusted facts are preserved and
whether conflicts are surfaced rather than silently overridden.

\paragraph{Attack Success Rate (ASR).}
The fraction of attack cases that induce an integrity failure (e.g., omission or contradiction of a trusted allergy).

\paragraph{Trusted Correct Rate.}
The fraction of cases where all required trusted facts are correctly reflected in the output.

\paragraph{Parse Error Rate.}
The fraction of generations that fail to satisfy the required output format (when a structured format is used).

\section{Results}

We report results for the allergy integrity scenario, which serves as a representative safety-critical task where silent omission or negation constitutes a serious trust violation. All results are computed over 1{,}000 attack cases per model unless otherwise noted.

\subsection{Baseline Vulnerability to EHR-Mediated Indirect Prompt Injection}

% AI-SUGGESTION: Insert baseline results table before first reference so \ref{tab:attack_baseline} resolves.
\begin{table}[t]
\centering
\caption{Baseline attack success rates for allergy omission (1{,}000 cases per model).}
\label{tab:attack_baseline}
\begin{tabular}{lr}
\toprule
Model & Attack success rate \\
\midrule
claude\_sonnet45 & 0.659 \\
gpt41           & 1.000 \\
gpt51           & 0.989 \\
llama31\_groq    & 0.362 \\
llama33\_groq    & 0.994 \\
qwen3\_bedrock   & 0.940 \\
\bottomrule
\end{tabular}
\end{table}

Table~\ref{tab:attack_baseline} summarizes attack success rates under the baseline EHR–LLM summarization pipeline without any trust-aware defenses. Across models, we observe consistently high vulnerability to EHR-mediated indirect prompt injection.

State-of-the-art general-purpose models exhibit particularly severe failures. GPT-4.1 shows a 100\% attack success rate, while GPT-5.1 and Claude Sonnet 4.5 exceed 98\% and 65\% respectively. Qwen3 deployed via Bedrock also demonstrates a high attack success rate of 94\%. These failures occur despite the attacks being embedded solely in patient-reported messages and remaining clinically plausible.

Open-weight models exhibit more variability. LLaMA~3.1 shows lower susceptibility (36.2\%), while LLaMA~3.3 is highly vulnerable (99.4\%). This suggests that vulnerability is not monotonic with model scale or perceived capability, and instead depends on how strongly the model prioritizes instruction-following over data integrity when faced with mixed-trust context.

Overall, these results demonstrate that EHR-mediated indirect prompt injection is a \emph{systemic vulnerability} affecting both proprietary and open models, rather than a flaw isolated to a specific architecture or training regime.

\subsection{Effectiveness of Trust-Aware Summarization Defense}

% AI-SUGGESTION: Insert defense results table before first reference so \ref{tab:defense_results} resolves.
\begin{table}[t]
\centering
\caption{Defense results for allergy omission (1{,}000 cases per model). Lower attack success is better; higher defense success is better.}
\label{tab:defense_results}
\begin{tabular}{lrrrr}
\toprule
Model & Clean trusted correct & Attack success & Defense success & Clean parse error \\
\midrule
claude45        & 0.756 & 0.245 & 0.755 & 0.244 \\
gpt41           & 0.995 & 0.000 & 1.000 & 0.000 \\
gpt51           & 0.987 & 0.005 & 0.995 & 0.000 \\
llama31\_groq    & 0.976 & 0.018 & 0.982 & 0.000 \\
llama33\_groq    & 0.973 & 0.034 & 0.966 & 0.000 \\
qwen3\_groq      & 0.994 & 0.008 & 0.992 & 0.000 \\
\bottomrule
\end{tabular}
\end{table}

We next evaluate the proposed trust-aware summarization defense, which separates trusted structured chart facts from untrusted patient-reported text and enforces integrity preservation during generation. Results are shown in Table~\ref{tab:defense_results}.

Across all evaluated models, the defense substantially reduces attack success rates. For GPT-4.1, GPT-5.1, and Qwen3, attack success rates drop from above 94\% to below 1\%, effectively eliminating silent allergy omission. Similar improvements are observed for both LLaMA variants, with post-defense attack success rates below 4\%.

Claude Sonnet 4.5 represents the most challenging case: while the defense reduces attack success from 65.9\% to 24.5\%, its performance is constrained by a relatively high output parsing error rate. Importantly, even in this case, the defense improves the trusted fact preservation rate from baseline levels and prevents the majority of silent integrity violations.

\subsection{Clean-Case Performance and Trade-offs}

A key requirement for clinical deployment is that defenses do not degrade performance under non-adversarial conditions. Across all models, clean trusted-fact correctness remains high after applying the defense, with most models maintaining correctness above 97\%. No model exhibits a meaningful drop in clean-case accuracy, indicating that the defense does not trade correctness for robustness.

Parsing error rates remain unchanged between clean and attack conditions, suggesting that observed improvements are attributable to trust-aware handling rather than artifacts of output filtering.

\subsection{Summary of Findings}

Taken together, the results demonstrate three key findings:
\begin{itemize}
    \item EHR-mediated indirect prompt injection reliably induces silent integrity violations across a wide range of LLMs, including state-of-the-art proprietary models.
    \item The vulnerability arises even when the attacker controls only patient-reported messages and does not require explicit jailbreak-style instructions.
    \item A lightweight, system-level trust-aware summarization strategy dramatically reduces attack success rates while preserving clean-case performance.
\end{itemize}

These findings reinforce that trustworthiness failures in medical LLM deployments stem primarily from mixed-trust context handling at the system level, rather than from insufficient model capability alone.

\section{Discussion}

\subsection{Why Do Strong Models Fail Under Mixed-Trust Context?}

A striking result of our study is that state-of-the-art general-purpose models, including GPT-4.1 and GPT-5.1, exhibit near-total failure under EHR-mediated indirect prompt injection. This failure is not due to a lack of factual knowledge or language understanding, but rather to how these models prioritize instruction-following behavior when confronted with mixed-trust inputs.

Patient-reported messages often contain language that implicitly frames preferences or instructions (e.g., requests for brevity, assertions about resolved conditions). When these messages are concatenated with structured EHR facts in a single context window, the model lacks an explicit mechanism to distinguish authoritative data from unverified text. As a result, even advanced models may incorrectly treat patient messages as higher-priority guidance, leading to silent omission or negation of safety-critical facts. This suggests that improved model capability alone does not resolve the problem; instead, explicit trust handling is required at the system level.

\subsection{System-Level Vulnerability Rather Than Model-Specific Flaws}

Our results show that vulnerability to EHR-mediated indirect prompt injection is widespread across model families, deployment modalities, and capability tiers. Both proprietary and open-weight models are affected, and susceptibility varies non-monotonically with perceived model strength. For example, LLaMA~3.1 exhibits lower attack success than LLaMA~3.3, while GPT-4.1 fails completely under attack.

This pattern indicates that the root cause lies in the application pipeline rather than in individual model architectures. Specifically, the assumption that an LLM will naturally preserve the integrity of structured chart facts when presented alongside untrusted text is unfounded. Without explicit trust boundaries, the summarization pipeline itself becomes an attack surface. From a systems perspective, this reframes prompt injection from a user-interface problem into a data provenance and integrity problem.

\subsection{Effectiveness and Limits of Trust-Aware Summarization}

The proposed trust-aware summarization defense substantially reduces attack success rates across all evaluated models, demonstrating that modest pipeline changes can dramatically improve trustworthiness. By separating trusted and untrusted sources and enforcing integrity preservation, the defense prevents silent overrides without requiring model retraining or access to internal weights.

However, the defense does not eliminate all failure modes. In particular, models with higher parsing error rates, such as Claude Sonnet~4.5 in our experiments, still exhibit residual vulnerability. This highlights an important limitation: system-level defenses depend on reliable output structure and may be constrained by model compliance with formatting or schema requirements. As such, trust-aware summarization should be viewed as a necessary but not sufficient component of a robust clinical deployment.

\subsection{Implications for Clinical Deployment}

From a clinical perspective, silent integrity violations are especially problematic because they may go unnoticed by clinicians who rely on summaries for rapid chart review. The fact that such violations can be triggered by benign patient messages underscores the practical risk of deploying LLM summarization tools without explicit trust safeguards.

Our findings suggest that clinical LLM deployments should treat patient-reported text as inherently untrusted and design pipelines that preserve the primacy of structured chart facts. Importantly, this does not require rejecting or ignoring patient input; rather, it requires surfacing conflicts explicitly and preventing unverified text from silently overriding verified records. This design principle aligns with existing clinical workflows, in which patient reports are routinely reviewed and validated by clinicians before updating the chart.

\subsection{Broader Implications for LLM Trustworthiness}

Although our study focuses on EHR summarization, the underlying issue extends to any LLM application that aggregates heterogeneous inputs with different trust levels. Knowledge bases, user-generated content, and external documents may all act as carriers of indirect prompt injection when combined with trusted data sources. Consequently, mixed-trust context handling should be considered a first-class concern in LLM system design.

Future work may explore integrating trust-aware summarization with complementary techniques such as automated verification, provenance tracking, or uncertainty-aware generation. More broadly, our results support the view that trustworthiness is a system property emerging from the interaction between models, data, and application logic, rather than a characteristic of the model alone.

% AI-SUGGESTION: Responsible/ethics framing for a clinical-security benchmark paper.
\subsection{Responsible Use, Ethics, and Limitations}
This work is intended to improve the safety of clinical LLM deployments by identifying a realistic systems vulnerability and evaluating lightweight mitigations. Our benchmark uses synthetic patient records and does not involve real patient data. Nonetheless, we recognize that publishing attack methodologies can enable misuse. To reduce misuse risk, we focus on clinically grounded but non-actionable prompt patterns and emphasize system-level mitigations that are implementable without model access. We do not claim that LLM-generated summaries are suitable for autonomous clinical decision-making; outputs should be treated as assistive drafts requiring clinician review. Finally, our defenses rely on structured output and automated checks; residual failures may occur due to formatting noncompliance or incomplete integrity rules, motivating future work on stronger provenance enforcement and comprehensive safety evaluation across additional clinical tasks.

\section{Conclusion}

This work investigates a previously underexplored trustworthiness failure in EHR-integrated large language models: silent integrity violations caused by EHR-mediated indirect prompt injection. Using a realistic clinical summarization setting, we show that patient-reported messages—when treated as untrusted but concatenated with authoritative chart data—can systematically induce omission or negation of safety-critical facts across a wide range of state-of-the-art LLMs.

Our results demonstrate that this vulnerability is not limited to weaker or open-weight models, but affects leading proprietary systems as well. The failure arises from a system-level assumption that LLMs can implicitly infer trust boundaries within mixed-context inputs, an assumption that our experiments show to be unreliable. Consequently, improvements in model capability alone are insufficient to address this class of failures.

To mitigate this issue, we propose a trust-aware summarization strategy that enforces separation between trusted structured EHR facts and untrusted patient-reported text during generation. Without requiring model retraining or access to internal parameters, this approach substantially reduces attack success rates while preserving clean-case performance. These findings suggest that explicit trust handling at the application level is a practical and effective path toward safer clinical LLM deployments.

More broadly, this work highlights the importance of viewing LLM trustworthiness as a property of the surrounding system rather than of the model in isolation. As LLMs continue to be integrated into safety-critical domains, future research should prioritize principled mechanisms for handling mixed-trust inputs, surfacing conflicts transparently, and aligning model behavior with domain-specific integrity requirements.

% =====================================================
\bibliographystyle{named}
\bibliography{ijcai26}

\end{document}
